% =====================================================================
%  第 20 章  重夸克物理：B 介子衰变与重夸克偶素            （ch20_heavy.tex，writer W23）
%  约定：本书约定（WRITER_SPEC_v2 §3）；n_A=(1^+,0^-,0_T)、n_B=(0^+,1^-,0_T)（不用 SCET 的 n、nbar）；
%  D_mu = d_mu - i g A_mu^a t^a；D = 4-2 eps，S_eps=(4pi)^eps e^{-eps gamma_E}；
%  BBL 与 Manohar--Wise 用 D = d + i g A（g -> -g）。验证脚本：writing/work/W23/s1_bdpi.py … s8_review.py（日志在 logs/）。
% =====================================================================
\chapter{重夸克物理：$B$ 介子衰变与重夸克偶素}\label{ch:heavy}

重夸克质量 $m_Q\gg\Lqcd$ 本身就是硬标度：$b$ 夸克衰变使末态轻介子以近光速运动，出现与硬遍举过程（\cref{ch:excl}）相同的共线结构；重夸克偶素中的相对速度 $v\ll1$ 提供第二个展开参数。本章只有一条严格结果——电磁衰变匹配系数的树图值与单圈值（[R, $R_{\mathrm{PT}}$]）；其余结果的地位从全阶论证到猜想不等，见下框。贯穿全章的一点是：衰变率是固定物理点上的 Minkowski 量，把这些结果用于固定运动学点上的强子衰变率还需要局域夸克—强子对偶（\cref{conj:ee-duality}(b)，[C]），它不在任何微扰证明的范围之内。

\begin{tcolorbox}[enhanced,breakable,title=本章证明了什么／没有证明什么,fonttitle=\bfseries,
  colback=gray!3!white,colframe=gray!55!black,boxrule=0.6pt,arc=2pt,
  left=5pt,right=5pt,top=3pt,bottom=3pt,before skip=6pt,after skip=6pt]
\textbf{证明了：}\cref{prop:heavy-nrqcd-matching}（\tagR，$R_{\mathrm{PT}}$：电磁衰变匹配系数的树图值与对 $\varepsilon$ 精确的单圈 $c_v$）。\par
\textbf{给出全阶论证：}\cref{thm:heavy-bdpi,prop:heavy-btopi-ff}（\tagA；瓶颈：\cref{rem:scet-assumptions}的 SCET 前提，$B\to D\pi$ 另需端点压低假设）；\cref{thm:heavy-inclusive-b}（\tagA；瓶颈：微扰 OPE \cref{thm:dis-ope}，端点区另需 SCET 前提）；\cref{thm:heavy-nrqcd-decay}（\tagA；瓶颈：拓扑因子化的全阶形式）；\cref{prop:heavy-quarkonium-fragmentation}（\tagA；瓶颈：强子对撞中有色末态的 Glauber 相消）。强子层次均另需\cref{asm:intro-bridge}。\par
\textbf{没有证明（或已知不成立）：}两个轻介子末态领头幂次因子化的全阶形式（\cref{prop:heavy-qcdf-lp}，\tagP）；其幂次修正的标准共线因子化——显式计算给出发散的卷积（\cref{prop:heavy-qcdf-endpoint}，\tagV）；产生的 NRQCD 因子化（\cref{conj:heavy-nrqcd-production}，\tagC）；固定运动学点上所需的局域对偶（\cref{conj:ee-duality}，\tagC）。
\end{tcolorbox}

\paragraph{记号。}
本章用 $n_A=(1^+,0^-,\mathbf 0_T)$、$n_B=(0^+,1^-,\mathbf 0_T)$，引用 SCET 文献（\cref{ch:scet}）时作换算 $n\to\sqrt2\,n_A$、$\bar n\to\sqrt2\,n_B$。被发射的轻介子或喷注总取为沿 $+z$ 运动，$p\simeq p^+n_A$：其组分经 $n_A\cdot A_s=A_s^-$ 与软胶子耦合，与之共线的胶子求和成沿 $n_B$ 的 Wilson 线。$\lambda=\Lqcd/Q$，$Q$ 为硬标度。重夸克有效理论（HQET）的场为 $h_v=\ee^{\ii m_Qv\cdot x}\frac{1+\sla v}{2}Q$，$v$ 为强子的四速度。BBL\cite{BodwinBraatenLepage1995}与 Manohar--Wise\cite{ManoharWise1994}用 $D_\mu=\partial_\mu+\ii gA_\mu$，其公式已按 $g\to-g$ 换算到\cref{def:prelim-qcd}的约定；含场强的式子一律用协变导数的对易子写出。
